\subsubsection{缓增分布}\label{subsec:tempered-distributions}

可以定义不同的试验函数空间, 从而定义相应的分布空间. 特别地, 对 $\Omega=\mathbb{R}^m$, 自然可以将紧支集要求替换为在无穷远处快速衰减的要求. 这引出以下定义.

\begin{Definition}
\label{def:rapidly-decreasing-test-space}
设 $\mathcal{S}(\mathbb{R}^m)$ 为 $\mathbb{R}^m$ 上所有属于 $C^\infty$ 类, 且对每个 $k\in\mathbb{N}$ 和每个多重指标 $\alpha$ 都使 $|\bm{x}|^k|D^\alpha\phi(\bm{x})|$ 有界的复值函数构成的空间. 若 $\phi_n$ 的所有阶导数均一致收敛到 $\phi$ 的相应导数, 且界 $|\bm{x}|^k|D^\alpha\phi_n(\bm{x})|\leq C_{k\alpha}$ 中的常数 $C_{k\alpha}$ 可不依赖于 $n$ 来选取, 则称 $\mathcal{S}(\mathbb{R}^m)$ 中的序列 $\phi_n$ 收敛到 $\phi$.
\end{Definition}

显然, $\mathcal{D}(\mathbb{R}^m)$ 是 $\mathcal{S}(\mathbb{R}^m)$ 的子空间. 此外, $\mathcal{D}(\mathbb{R}^m)$ 在 $\mathcal{S}(\mathbb{R}^m)$ 中稠密. 为说明这一点, 设 $e(\bm{x})$ 为 $C^\infty$ 函数, 在单位球内等于 1, 在半径为 2 的球外为零. 令 $e_n(\bm{x})=e(\bm{x}/n)$. 那么对任意 $f\in\mathcal{S}(\mathbb{R}^m)$, 有 $f=\lim_{n\to\infty}fe_n$.

现在将缓增分布定义为 $\mathcal{S}$ 上的连续线性泛函.

\begin{Definition}
\label{def:tempered-distribution}
$\mathbb{R}^m$ 上的\textbf{缓增分布}是从 $\mathcal{S}(\mathbb{R}^m)$ 到 $\mathbb{C}$ 的线性映射 $\phi\mapsto(f,\phi)$, 它具有如下连续性: 若在 $\mathcal{S}(\mathbb{R}^m)$ 中 $\phi_n\to\phi$, 则 $(f,\phi_n)\to(f,\phi)$. 所有缓增分布构成的集合记为 $\mathcal{S}^{\prime}(\mathbb{R}^m)$. 若对每个 $\phi\in\mathcal{S}(\mathbb{R}^m)$ 都有 $(f_n,\phi)\to(f,\phi)$, 则称在 $\mathcal{S}^{\prime}(\mathbb{R}^m)$ 中 $f_n\to f$.
\end{Definition}

显然, 每个缓增分布通过限制都定义一个分布. 此外, 若两个缓增分布作为 $\mathcal{D}^{\prime}(\mathbb{R}^m)$ 的元素相等, 则它们作为 $\mathcal{S}^{\prime}(\mathbb{R}^m)$ 的元素也相等; 这来自 $\mathcal{D}(\mathbb{R}^m)$ 在 $\mathcal{S}(\mathbb{R}^m)$ 中稠密这一事实. 因而 $\mathcal{S}^{\prime}(\mathbb{R}^m)$ 是 $\mathcal{D}^{\prime}(\mathbb{R}^m)$ 的线性子空间. 此外, $\mathcal{S}^{\prime}(\mathbb{R}^m)$ 中的收敛显然蕴含 $\mathcal{D}^{\prime}(\mathbb{R}^m)$ 中的收敛.
